Multimodal sentiment analysis (MSA) integrates language, acoustic, and visual cues,
yet many fusion architectures allow low-level or noisy auxiliary signals to
shortcut task-relevant structure, hurting calibration and robustness.
We present \textbf{PRISM}, a \emph{primary-centric} model that compresses each
modality with variational information-bottleneck (IB) encoders, exposes per-token
confidence from IB uncertainty, and fuses modalities through confidence-biased
multi-head attention together with learned adaptive gates.
A differentiable \textbf{MSelector} picks a per-sample primary modality using
Gumbel--Softmax routing; stage-two training adds selective KL supervision that
aligns routing weights with per-modality predictive errors.
We train with a two-stage curriculum, AdamW optimization, optional exponential
moving averages (EMA) for evaluation, and checkpoint selection on development
MAE (regression benchmarks).
On CMU-MOSI and CMU-MOSEI, a representative configuration achieves
87.94\% binary accuracy and 0.605 MAE on MOSI, and 87.21\% / 0.599 MAE on MOSEI
(complete-modality setting).
On CH-SIMS v2, an Optuna tier-1 search (80 trials) reaches dev-selected test
MAE $\approx$0.319 with Acc-2 $\approx$78\%.
We discuss redundancy between denoising front-ends and IB compression, and
outline token-level stochastic extensions suggested by concurrent CyIN-style
designs \citep{cyin2025align}.
